\chapter{Proceso de Ingesta ETL Detallado}\label{anexo-d.-proceso-de-ingesta-etl-detallado}

El proceso ETL (\textit{geotiff\_to\_tiledb\_sparse.py}) transforma los ficheros GeoTIFF generados por Triton en un array TileDB sparse listo para consulta. Este anexo describe el proceso con el nivel de detalle necesario para reproducirlo o adaptarlo a otro simulador.

\section{Parámetros de configuración}\label{d.1-parametros-de-configuracion}

Los tres parámetros que controlan el comportamiento del ETL se definen como constantes al inicio del script:

\begin{table}[ht]
\centering
\begin{tabularx}{\textwidth}{>{\raggedright\arraybackslash}p{4.5cm}L>{\raggedright\arraybackslash}p{2.5cm}}
\toprule
\textbf{Parámetro} & \textbf{Descripción} & \textbf{Valor} \\
\midrule
\texttt{DEPTH\_THRESHOLD} & Umbral de celda húmeda: celdas con H por debajo de este valor se descartan y no se almacenan en TileDB & 0,01 m \\
\texttt{CHUNK\_ROWS} & Número de filas que se leen simultáneamente de cada GeoTIFF. Controla el consumo máximo de RAM por franja & 5\,000 filas \\
tile (espacial) & Tamaño del tile en las dimensiones fila y columna del array TileDB & 256 $\times$ 256 \\
\bottomrule
\end{tabularx}
\caption{Parámetros configurables del proceso ETL.}
\label{tab:etl-params}
\end{table}

\section{Esquema del array TileDB}\label{d.2-esquema-del-array}

El script llama a \texttt{create\_sparse\_array()} para crear el array con la siguiente configuración:

\begin{itemize}
\item \textbf{Modo:} disperso (\texttt{sparse=True}), sin duplicados (\texttt{allows\_duplicates=False}).
\item \textbf{Orden de celdas y tiles:} \texttt{row-major} en ambos casos.
\item \textbf{Dimensiones:} \texttt{time} (int32, rango 0--19, tile=16), \texttt{row} (int32, rango 0--34\,199, tile=256), \texttt{col} (int32, rango 0--23\,039, tile=256).
\item \textbf{Atributos:} H, QX, QY, MH (todos float32), con el filtro ByteShuffle + ZSTD-9 aplicado individualmente a cada atributo.
\item \textbf{Filtro de compresión:} \texttt{tiledb.ByteShuffleFilter()} seguido de \texttt{tiledb.ZstdFilter(level=9)}, elegidos tras el benchmark del \hyperref[anexo-b.-tablas-de-resultados-completas]{Anexo~B} (reducción adicional del 8,9\,\% frente a ZSTD-1 sin penalización en lectura).
\end{itemize}

Si el array ya existe en disco, el script lo elimina antes de recrearlo para evitar la acumulación de fragmentos de ejecuciones anteriores.

\section{Bucle de procesamiento por paso temporal}\label{d.3-bucle-por-paso}

El cuerpo principal del ETL itera sobre los 20 pasos temporales en orden cronológico. Para cada paso \texttt{t}:

\begin{enumerate}
\item \textbf{Apertura de GeoTIFFs.} Se abren simultáneamente los cuatro ficheros correspondientes al paso (\texttt{H\_<step>.tif}, \texttt{QX\_<step>.tif}, \texttt{QY\_<step>.tif}, \texttt{MH\_<step>.tif}) con la biblioteca rasterio.

\item \textbf{Lectura por franjas.} El dominio completo (34\,200 $\times$ 23\,040 píxeles) se recorre en franjas de \texttt{CHUNK\_ROWS} = 5\,000 filas. Cada franja se lee con \texttt{rasterio.windows.Window} y la lectura de las cuatro bandas se paraleliza con \texttt{ThreadPoolExecutor(max\_workers=4)}, reduciendo el tiempo de E/S respecto a la lectura secuencial.

\item \textbf{Filtrado de celdas húmedas.} Sobre la franja leída se aplica la máscara \texttt{H >= 0,01 m}. Las coordenadas locales de las celdas que pasan el filtro se convierten a coordenadas globales (sumando el desplazamiento de fila \texttt{row\_off}) y se acumulan en listas junto con los valores de los cuatro atributos.

\item \textbf{Escritura en TileDB.} Una vez procesadas todas las franjas del paso temporal, se realiza \textbf{una única escritura} con el total de celdas húmedas del paso. Esto crea exactamente un fragmento en el array TileDB, lo que preserva la separación temporal que hace eficiente el filtrado por tiempo (véase la Sección~\ref{fragmentos-y-consolidacion}).

\item \textbf{Progreso.} El script imprime por pantalla el número de celdas húmedas almacenadas y el tiempo empleado en cada paso.
\end{enumerate}

\section{Análisis de memoria}\label{d.4-analisis-de-memoria}

El diseño por franjas limita el consumo de RAM a dos contribuciones principales:

\begin{itemize}
\item \textbf{Buffers de lectura:} una franja de 5\,000 filas $\times$ 23\,040 columnas $\times$ 4 variables $\times$ 4 bytes/float32 $\approx$ 1,8 GB.
\item \textbf{Acumuladores de celdas húmedas:} listas de coordenadas y valores para el paso completo. En el caso más desfavorable (paso con mayor fracción húmeda), pueden alcanzar varios cientos de millones de entradas; para una fracción húmeda típica del 8,3\,\% esto supone $\sim$65 millones de celdas $\times$ 6 arrays (2 coords + 4 attrs) $\times$ 4 bytes $\approx$ 1,5 GB.
\end{itemize}

El consumo pico total se sitúa alrededor de los 3--4 GB de RAM, lo que hace recomendable disponer de al menos 8 GB libres durante la ingesta.

\section{Metadatos almacenados}\label{d.5-metadatos}

Tras completar la escritura de los 20 fragmentos, el script almacena en el propio array TileDB los metadatos geoespaciales necesarios para interpretar los datos y para la interoperabilidad con herramientas SIG:

\begin{table}[ht]
\centering
\begin{tabularx}{\textwidth}{>{\raggedright\arraybackslash}p{4.5cm}L}
\toprule
\textbf{Clave} & \textbf{Contenido} \\
\midrule
\texttt{crs}        & Sistema de referencia de coordenadas (EPSG:32614, UTM zona 14N) \\
\texttt{transform}  & Transformación afín GeoTIFF (6 coeficientes) en JSON \\
\texttt{bounds}     & Extensión espacial en coordenadas del CRS (JSON) \\
\texttt{width} / \texttt{height} & Dimensiones del ráster original \\
\texttt{time\_steps} & Lista de identificadores de paso temporal en JSON \\
\texttt{source\_dataset} & Alias del \textit{dataset} (p.~ej.\ \texttt{datos1}) \\
\texttt{depth\_threshold\_m} & Umbral de celda húmeda usado en la ingesta \\
\texttt{storage\_mode} & Cadena descriptiva de la configuración de compresión \\
\texttt{zero\_rule} & Convención: celda ausente equivale a H=QX=QY=MH=0 \\
\bottomrule
\end{tabularx}
\caption{Metadatos geoespaciales almacenados en el array TileDB tras la ingesta.}
\label{tab:etl-metadata}
\end{table}

Estos metadatos permiten al motor de consultas reconstruir el georeferenciado sin necesidad de conservar los GeoTIFF originales, lo que hace posible desplegar el sistema en un servidor donde solo se transfieren los arrays TileDB (17--21 GB por escenario) y no los datos brutos (~235 GB). La clave \texttt{zero\_rule} garantiza la interoperabilidad: cualquier celda ausente en el array equivale a valor cero, convención necesaria para que la exportación a GeoTIFF produzca rásters completos con nodata correctamente asignado.

\section{Ejecución e instrucciones de uso}\label{d.6-ejecucion}

El proceso de ingesta se compone de tres fases: verificación de requisitos, ejecución del script y comprobación de la salida generada.

\subsection{Requisitos previos}\label{d.6.1-requisitos}

\begin{itemize}
\item Los 80 ficheros GeoTIFF del escenario deben estar disponibles en el directorio fuente apuntado por \texttt{TRITON\_GTIFF\_DIR/<dataset>}.
\item La variable de entorno \texttt{TRITON\_BASE\_URI} debe apuntar al directorio donde se crearán los arrays TileDB. Si no se define, se utiliza el valor por defecto configurado en \texttt{config.py}.
\item El entorno Python debe incluir las dependencias: \texttt{tiledb>=0.30}, \texttt{rasterio>=1.3}, \texttt{numpy>=1.24}.
\end{itemize}

Una vez verificados estos tres requisitos, el script puede ejecutarse con el siguiente comando.

\subsection{Comando de ejecución}\label{d.6.2-comando}

\begin{verbatim}
cd spillway/
python geotiff_to_tiledb_sparse.py --dataset datos1
\end{verbatim}

Si las variables de entorno no están exportadas en la sesión, pueden pasarse en línea:

\begin{verbatim}
TRITON_BASE_URI=/ruta/triton_results \
TRITON_GTIFF_DIR=/ruta/datos \
    python geotiff_to_tiledb_sparse.py --dataset datos1
\end{verbatim}

\subsection{Salida esperada}\label{d.6.3-salida}

Al finalizar la ingesta el script imprime un resumen similar al siguiente (valores reales de datos1):

\begin{verbatim}
Dataset: datos1
  GTIFF_DIR:  $TRITON_GTIFF_DIR/datos1
  TILEDB_URI: .../output_10_..._19800728_19800801
Pasos encontrados: 20
Resolución: 23040 x 34200 px  |  Franja: 5000 filas
  [ 1/20] paso 00_00 — 56,862,388 celdas mojadas  (23.4s)
  [ 2/20] paso 06_00 — 57,340,112 celdas mojadas  (21.8s)
  ...
  [20/20] paso 114_00 — 74,123,946 celdas mojadas (24.1s)

Array TileDB sparse creado en:    .../output_10_...
Pasos temporales:                  20
Total celdas almacenadas:          1,311,695,098
Umbral aplicado:                   H >= 0.01 m
Tiempo total:                      454.3s
\end{verbatim}

El tiempo total varía entre 7 y 15 minutos según el \textit{dataset} (el tamaño de \texttt{datos29}, el más grande con 21 GB, requiere aproximadamente el doble de tiempo). El número de fragmentos resultante debe ser exactamente 20; puede verificarse con:

\begin{verbatim}
python - <<'EOF'
import tiledb, sys
uri = sys.argv[1]
fi = tiledb.array_fragments(uri)
print(f"Fragmentos: {len(fi)}")  # debe ser 20
EOF
\end{verbatim}

Un número de fragmentos distinto de 20 (en particular, 1) indica que se aplicó consolidación no deseada, lo que degrada el rendimiento de las consultas multi-paso (véase la Sección~\ref{fragmentos-y-consolidacion}).
